\begin{textsample}{POS Dim 1 – vem\_ed\_32 – Score 24.00 – vem\_ed\_32/t172.txt}  \label{ex:f1_pos_033}
Intercâmbios Virtuais como Ponte de Empatia e Inovação no Ensino Artigo de Opinião Após \textbf{ser} convidada para participar de uma mesa-redonda no \textbf{início} do mês de outubro deste ano, pedi \textbf{que} alguns alunos gravassem vídeos relatando suas experiências em Projetos Colaborativos Internacionais (PCIs / CGESG), conhecidos mundialmente como Collaborative Online International Learning (COIL).
Como professora de inglês na Fatec Guaratinguetá, já participei de PCIs com instituições dos Estados Unidos, Armênia, Filipinas, México e Israel.
Essas parcerias sempre me pareceram uma oportunidade pedagógica, mas ouvir os estudantes falando me fez perceber o \textbf{impacto} \textbf{humano} e transformador desses projetos.
Muitos alunos relataram terem se sentido mais confiantes depois de interagir com colegas de outros países.
Alguns disseram \textbf{que}, graças ao PCI, se prepararam para entrevistas de emprego em inglês e conseguiram oportunidades profissio- nais.
Outros \textbf{contaram} \textbf{que} o \textbf{intercâmbio} virtual reacendeu seu interesse pela língua inglesa: não mais apenas um \textbf{conteúdo} da escola, mas algo com significado real em suas vidas.
Um dos depoimentos mais marcantes foi de um aluno \textbf{que} lutava contra a depressão.
Ele disse \textbf{que}, depois de participar de um PCI, redescobriu sua paixão pelo inglês, \textbf{criou} um canal no YouTube e encontrou uma nova razão para continuar estudando.
Para mim, esse tipo de \textbf{transformação} mostra \textbf{que} o PCI / COIL \textbf{é} mais do \textbf{que} ensino: \textbf{é} uma força de \textbf{mudança} social.
No \textbf{nível} docente, também experimento um profundo crescimento.
Planejar cursos com parceiros internacionais exige empatia, cooperação e adaptação — e nos deixa mais abertos às \textbf{realidades} dos outros.
Mesmo com a distância física, percebo \textbf{que} muitos desafios educacionais \textbf{que} enfrentamos (burocracia, limita- ção tecnológica, engajamento) \textbf{são} \textbf{compartilhados} por professores ao redor do mundo.
Isso \textbf{fortalece} a sensação de comunidade \textbf{profissional} global.
Além disso, os projetos ajudam a cultivar cidadania global.
Os alunos percebem \textbf{que} suas vozes têm relevância além de sua região, \textbf{que} podem dialogar com outras culturas e refletir sobre suas próprias identidades.
Esses \textbf{intercâmbios} os incentivam a pensar globalmente e agir localmente.
Re- conheço \textbf{que} há obstáculos reais: fusos \textbf{horários}, infraestrutura desigual e restrições tecnológicas.
Mas acredito \textbf{que} esses desafios se transformam em oportunidades.
Ao buscar ajustes, \textbf{somos} obrigados a repensar métodos, \textbf{cronogramas} e abordagens, e isso \textbf{gera} inovação e resiliência.
Por fim, afirmo \textbf{que} os projetos COIL merecem \textbf{apoio} \textbf{institucional} mais robusto: \textbf{é} preciso tempo para planejamento, recursos tecnológicos, capacitação docente.
Esse investimento não \textbf{é} apenas estratégico: \textbf{é} humanitário.
Quando estudantes relatam \textbf{que} o COIL lhes deu confiança, propósito ou alegria, percebemos \textbf{que} a internacionalização virtual não serve apenas para “fazer \textbf{intercâmbio}” — ela pode mudar vidas.
E, para nós, professores, \textbf{é} a confirmação de \textbf{que} ensinar \textbf{é}, antes de tudo, ligar \textbf{pontes} \textbf{humanas}. \textbf{Referências} HACKETT, S.; DAWSON, M.; JANSSEN, J.; VAN TARTWIJK, J.
Defining Collaborative Online International Learning (COIL) and distinguishing it from Virtual Exchange.
TechTrends, v. 68, n. 6, p. 1078-1094, 2024.
Disponível em: https://doi.org/10.1007/s11528-024-01000-w.
Acesso em: 28 set. 2025.
HACKETT, S.; JANSSEN, J.; BEACH, P.; PERREAULT, M.; BEELEN, J.; VAN TARTWIJK, J.
The effectiveness of Collaborative Online International Learning (COIL) on intercultural competence development in higher education.
International Journal of Educational Technology in Higher Education, v. 20, n. 1, art. 5, 2023.
Disponível em: https://doi.org/10.1186/s41239-022-00373-3.
Acesso em: 28 set. 2025.
RAJAGOPALAN, Kanavillil.
Vencer barreiras e emergir das adversidades com pleno êxito, sempre com o pé no chão.
In: LIMA, D.
C. de (Org.).
Inglês em escolas públicas não funciona?
Uma \textbf{questão}, múltiplos olhares.
São Paulo: Parábola, 2011, p. 55–65.

% matched lemmas: apoio, compartilhar, contar, conteúdo, criar, cronograma, fortalecer, gerar, horário, humano, impacto, institucional, intercâmbio, início, mudança, nível, ponte, profissional, que, questão, realidade, referência, ser, transformação
\end{textsample}
